% J3.2: del código suelto al paquete probado, documentado e integrado.
\documentclass[tikz,border=4pt]{standalone}
\usepackage{resonancia}
\begin{document}
\begin{tikzpicture}
  \node[cajaAzul, minimum width=2.4cm] (a) at (0,0) {\textbf{función}\\[-1pt]{\footnotesize en una página}};
  \node[cajaVerde, minimum width=2.4cm] (b) at (3.4,0) {\textbf{módulo}\\[-1pt]{\footnotesize \texttt{module}, \texttt{export}}};
  \node[cajaNaranja, minimum width=2.4cm] (c) at (6.8,0) {\textbf{paquete}\\[-1pt]{\footnotesize \texttt{Project.toml}}};
  \node[caja, draw=rVioleta, fill=rVioleta!8, minimum width=2.4cm] (d) at (10.2,0) {\textbf{pruebas}\\[-1pt]{\footnotesize \texttt{@testset}}};
  \node[caja, draw=rGris, fill=rGrisClaro!50, minimum width=2.4cm] (e) at (13.6,0) {\textbf{docs y CI}\\[-1pt]{\footnotesize Documenter, Actions}};
  \foreach \x/\y in {a/b, b/c, c/d, d/e} \draw[flecha] (\x) -- (\y);
  \node[etiqueta, color=rGris, anchor=north] at (6.8,-0.9) {cuando una función se reutiliza en varias páginas, sube un escalón};
\end{tikzpicture}
\end{document}
